\section{Použité nástroje}
Při implementaci bylo využito několik nástrojů, které zefektivňují práci při~vývoji softwaru. Pro hostování repozitáře se zdrojovým kódem byla využita platforma GitHub. Součástí repozitáře je rovněž automatizovaný bezpečnostní nástroj Dependabot, který byl aktivován za účelem průběžného monitorování závislostí projektu. Nástroj pravidelně kontroluje obsah souborů \texttt{composer.lock} a \texttt{package-lock.json} a v případě detekce známé bezpečnostní zranitelnosti v některé z použitých knihoven automaticky vytváří pull request s aktualizovanou verzí balíčku. Tím je zajištěna vysoká úroveň bezpečnosti aplikace bez nutnosti manuální kontroly všech externích komponent.

Jako vývojové prostředí byl při implementaci využit nástroj PhpStorm od~společnosti JetBrains. Tento editor je v komunitě PHP vývojářů široce využíván díky své robustní podpoře refaktoringu, pokročilé statické analýze kódu a integrovaným nástrojům pro ladění a testování. \cite{php-jetbrains} Pro zefektivnění vývoje nad frameworkem Laravel bylo prostředí rozšířeno o plugin Laravel Idea. Toto rozšíření poskytuje hlubší porozumění struktuře projektu, nabízí inteligentní doplňování kódu specifické pro Laravel a usnadňuje práci s Eloquent modely.

Pro zajištění konzistentního běhového prostředí byla využita technologie Docker. Aplikace je spolu se všemi svými závislostmi, jako je například webový server, zapouzdřena do izolovaných kontejnerů. Díky tomu bylo eliminováno riziko nekompatibility mezi lokálním prostředím a produkčním serverem. Definice infrastruktury v souboru \texttt{docker-compose.yml} navíc umožňuje rychlé sestavení aplikačních kontejnerů na jakémkoli zařízení s podporou Dockeru.
